% Chapter 18: Sensitivity Analysis for the Average Causal Effect
% A First Course in Causal Inference - Peng Ding
% Rewrite V2

\chapter{平均因果效应的灵敏度分析}\label{ch:18}

\section{引言：可忽略性失效时我们能做什么}\label{sec:18-intro}

在 Part III 中，我们详细讨论了四种基于可忽略性假设的因果推断方法：outcome regression（第 \ref{ch:10} 章）、逆倾向得分加权（第 \ref{ch:11} 章）、双重稳健估计（第 \ref{ch:12} 章）以及匹配（第 \ref{ch:15} 章）。可忽略性假设意味着在协变量 $X$ 条件下，治疗分配 $Z$ 与潜在结果相互独立：
\[
Z \perp \{Y(1), Y(0)\} \mid X .
\]
这个假设保证了我们可以从观测数据中识别平均因果效应 $\tau = \mathbb{E}\{Y(1) - Y(0)\}$。

然而，第 \ref{ch:17} 章已经指出，可忽略性假设无法从数据本身验证。当研究者无法控制所有混淆因素时，观测到的治疗与结果之间的关联可能是虚假的。

面对这一困境，统计学提供了两条出路：

\textbf{第一条路（乐观）}：假设可忽略性成立，利用观测数据直接估计因果效应。这是 Part III 的核心思想。

\textbf{第二条路（悲观）}：完全放弃从观测数据中推断反事实，转而报告因果效应在所有可能情况下的最坏情况边界。Manski (1990) 的界是这一思路的典型代表。

\textbf{第三条路（折中）}：既不盲目假设可忽略性，也不完全放弃推断，而是引入有限的灵敏度参数，量化未测混杂的强度对因果效应估计的影响程度。本章正是沿着这一思路展开。

第 \ref{ch:17} 章的 E-value 量化了"要解释掉观察到的关联，未测混杂需要有多强"。本章则将这一思想系统化，基于 outcome regression、IPW 和双重稳健估计，给出平均因果效应的完整灵敏度分析框架。

\section{Manski 型最坏情况界}\label{sec:18-manski-bounds}

\subsection{部分识别}\label{sec:18-partial-identification}

在讨论灵敏度分析之前，我们先介绍一个更弱但更基本的概念——\textbf{部分识别（partial identification）}。

\begin{Definition}[部分识别]\label{def:18-partial-id}
如果观测数据分布与参数 $\theta$ 的多个值相容，则称参数 $\theta$ 是\textbf{部分识别}的。
\end{Definition}

这与第 \ref{def:10-1} 章的定义形成对比：如果参数 $\theta$ 由观测数据分布唯一确定，则它是可识别的；否则，它只是部分识别的。

在可忽略性假设下，$\tau$ 是可识别的；但在没有可忽略性假设的情况下，$\tau$ 通常只是部分识别的。

\subsection{仅假设 outcomes 有界}\label{sec:18-bounded-outcomes}

我们仅假设 outcomes 有界：
\[
\underline{y} \leq Y_i(1), Y_i(0) \leq \overline{y} .
\]
对于二值结果，这等价于 $[\underline{y}, \overline{y}] = [0, 1]$。

在这个极弱假设下，我们可以导出 $E\{Y(1)\}$ 和 $E\{Y(0)\}$ 的最坏情况边界。考虑 $E\{Y(1)\}$ 的分解：
\[
E\{Y(1)\} = E\{Y(1) \mid Z=1\}\Pr(Z=1) + E\{Y(1) \mid Z=0\}\Pr(Z=0) .
\]

由于 $Y(1) \in [\underline{y}, \overline{y}]$，我们有
\[
E\{Y(1) \mid Z=0\} \in [\underline{y}, \overline{y}] .
\]

因此，$E\{Y(1)\}$ 的下界为
\[
E\{Y \mid Z=1\}\Pr(Z=1) + \underline{y}\Pr(Z=0) ,
\]
上界为
\[
E\{Y \mid Z=1\}\Pr(Z=1) + \overline{y}\Pr(Z=0) .
\]\footnote{推导见附录 \cref{sec:appendix-18-manski-bounds}}

同理，$E\{Y(0)\}$ 的下界为
\[
\underline{y}\Pr(Z=1) + E\{Y \mid Z=0\}\Pr(Z=0) ,
\]
上界为
\[
\overline{y}\Pr(Z=1) + E\{Y \mid Z=0\}\Pr(Z=0) .
\]

综合可得，平均因果效应 $\tau = E\{Y(1)\} - E\{Y(0)\}$ 的下界为
\[
E\{Y \mid Z=1\}\Pr(Z=1) + \underline{y}\Pr(Z=0) - \overline{y}\Pr(Z=1) - E\{Y \mid Z=0\}\Pr(Z=0) ,
\]
上界为
\[
E\{Y \mid Z=1\}\Pr(Z=1) + \overline{y}\Pr(Z=0) - \underline{y}\Pr(Z=1) - E\{Y \mid Z=0\}\Pr(Z=0) .
\]

这些边界的宽度为 $\overline{y} - \underline{y}$，比先验边界 $[\underline{y} - \overline{y}, \overline{y} - \underline{y}]$（宽度为 $2(\overline{y} - \underline{y})$）要窄，但通常仍然包含 0，因此缺乏信息量。

\textbf{历史注记}：Cochran (1953) 在缺失数据的调查中使用了最坏情况界思想，但因其结果往往过于保守而放弃。Manski (1990) 将这一思想引入因果推断，Manski (2003) 进一步在计量经济学模型中系统化了这一方法。

\section{平均因果效应的灵敏度分析}\label{sec:18-sensitivity-analysis}

上一节的 Manski 型界虽然理论优美，但实践中过于保守。本节引入灵敏度参数，提供一个更实用的折中方案。

\subsection{灵敏度参数}\label{sec:18-sensitivity-parameters}

定义两个灵敏度参数：
\begin{equation}
\varepsilon_1(X) = \frac{E\{Y(1) \mid Z=1, X\}}{E\{Y(1) \mid Z=0, X\}} , \quad
\varepsilon_0(X) = \frac{E\{Y(0) \mid Z=1, X\}}{E\{Y(0) \mid Z=0, X\}} .
\label{eq:18-sensitivity-params}
\end{equation}

直观上，$\varepsilon_1(X)$ 衡量的是在协变量 $X$ 条件下，治疗组对照组的潜在结果均值之比。当 $\varepsilon_1(X) = 1$ 时，治疗组和对照组的 $Y(1)$ 条件均值相等，说明没有未测混杂影响 $Y(1)$。

为简化，我们可以进一步假设灵敏度参数是常数（不依赖 $X$）：$\varepsilon_1(X) = \varepsilon_1$，$\varepsilon_0(X) = \varepsilon_0$。实践中，我们需要固定这些参数或在一个预先指定的范围内变化它们。

记 $\mu_1(X) = E(Y \mid Z=1, X)$ 和 $\mu_0(X) = E(Y \mid Z=0, X)$ 为观测结果在治疗组和对照组下的条件均值函数。

\subsection{识别公式}\label{sec:18-identification-formulas}

下面的定理给出了在已知灵敏度参数条件下，平均因果效应的识别公式。

\begin{Theorem}[识别公式 I：基于 outcome regression]\label{th:18-1}
在已知 $\varepsilon_1(X)$ 和 $\varepsilon_0(X)$ 的条件下，我们有
\[
E\{Y(1) \mid Z=0\} = E\left\{\frac{\mu_1(X)}{\varepsilon_1(X)} \mid Z=0\right\} ,
\]
\[
E\{Y(0) \mid Z=1\} = E\left\{\mu_0(X)\varepsilon_0(X) \mid Z=1\right\} .
\]

因此，
\begin{align}
\tau &= E\left[ZY + (1-Z)\frac{\mu_1(X)}{\varepsilon_1(X)}\right] - E\left[Z\mu_0(X)\varepsilon_0(X) + (1-Z)Y\right] \label{eq:18-ace-regression} \\
&= E\left[Z\mu_1(X) + (1-Z)\frac{\mu_1(X)}{\varepsilon_1(X)}\right] - E\left[Z\mu_0(X)\varepsilon_0(X) + (1-Z)\mu_0(X)\right] . \label{eq:18-ace-regression-alt}
\end{align}
\end{Theorem}

\textbf{定理 \ref{th:18-1} 的含义}：当 $\varepsilon_1(X) = \varepsilon_0(X) = 1$ 时，公式 \cref{eq:18-ace-regression} 简化为标准的 outcome regression 估计量。这验证了灵敏度参数是可忽略性假设的推广。\footnote{完整证明推导见附录 \cref{sec:appendix-18-theorem-1}}

基于 \cref{eq:18-ace-regression}，我们可以构造两类估计量：

\textbf{Predictive 估计量}（使用观测结果）：
\begin{align}
\hat{\tau}^{\mathrm{pred}} &= \left\{n^{-1}\sum_{i=1}^n Z_i Y_i + n^{-1}\sum_{i=1}^n (1-Z_i)\frac{\hat{\mu}_1(X_i)}{\varepsilon_1(X_i)}\right\} \label{eq:18-tau-pred} \\
&\quad - \left\{n^{-1}\sum_{i=1}^n Z_i\hat{\mu}_0(X_i)\varepsilon_0(X_i) + n^{-1}\sum_{i=1}^n (1-Z_i)Y_i\right\} . \\
\end{align}

\textbf{Projective 估计量}（使用拟合值）：
\begin{align}
\hat{\tau}^{\mathrm{proj}} &= \left\{n^{-1}\sum_{i=1}^n Z_i\hat{\mu}_1(X_i) + n^{-1}\sum_{i=1}^n (1-Z_i)\frac{\hat{\mu}_1(X_i)}{\varepsilon_1(X_i)}\right\} \label{eq:18-tau-proj} \\
&\quad - \left\{n^{-1}\sum_{i=1}^n Z_i\hat{\mu}_0(X_i)\varepsilon_0(X_i) + n^{-1}\sum_{i=1}^n (1-Z_i)\hat{\mu}_0(X_i)\right\} . \\
\end{align}

\textbf{术语说明}：``predictive'' 和 ``projective'' 的术语来自调查抽样文献（Firth and Bennett, 1998; Ding and Li, 2018）。前者使用观测结果，后者用拟合值替换。

下面的定理给出了基于逆概率加权的识别公式。

\begin{Theorem}[识别公式 II：基于 IPW]\label{th:18-2}
在已知 $\varepsilon_1(X)$ 和 $\varepsilon_0(X)$ 的条件下，
\[
E\{Y(1)\} = E\left\{w_1(X)\frac{Z}{e(X)}Y\right\} , \quad
E\{Y(0)\} = E\left\{w_0(X)\frac{1-Z}{1-e(X)}Y\right\} ,
\]
其中
\begin{equation}
w_1(X) = e(X) + \frac{1-e(X)}{\varepsilon_1(X)} , \quad
w_0(X) = e(X)\varepsilon_0(X) + 1 - e(X) .
\label{eq:18-ipw-weights}
\end{equation}
\end{Theorem}

\textbf{定理 \ref{th:18-2} 的含义}：经典 IPW 公式被两个额外的因子 $w_1(X)$ 和 $w_0(X)$ 修正，它们同时依赖于倾向得分和灵敏度参数。当 $\varepsilon_1(X) = \varepsilon_0(X) = 1$ 时，$w_1(X) = w_0(X) = 1$，公式退化为标准 IPW。\footnote{完整证明推导见附录 \cref{sec:appendix-18-theorem-2}}

基于 \cref{eq:18-ipw-weights}，可以构造 Horvitz-Thompson 型估计量：
\begin{align}
\hat{\tau}^{\mathrm{ht}} &= n^{-1}\sum_{i=1}^n \frac{\{\hat{e}(X_i)\varepsilon_1(X_i) + 1 - \hat{e}(X_i)\}Z_i Y_i}{\varepsilon_1(X_i)\hat{e}(X_i)} \label{eq:18-tau-ht} \\
&\quad - n^{-1}\sum_{i=1}^n \frac{\{\hat{e}(X_i)\varepsilon_0(X_i) + 1 - \hat{e}(X_i)\}(1-Z_i)Y_i}{1 - \hat{e}(X_i)} , \\
\end{align}
以及 Hajek 型估计量：
\begin{align}
\hat{\tau}^{\mathrm{haj}} &= \frac{\sum_{i=1}^n \frac{\{\hat{e}(X_i)\varepsilon_1(X_i) + 1 - \hat{e}(X_i)\}Z_i Y_i}{\varepsilon_1(X_i)\hat{e}(X_i)}}{\sum_{i=1}^n \frac{Z_i}{\hat{e}(X_i)}} \label{eq:18-tau-haj} \\
&\quad - \frac{\sum_{i=1}^n \frac{\{\hat{e}(X_i)\varepsilon_0(X_i) + 1 - \hat{e}(X_i)\}(1-Z_i)Y_i}{1 - \hat{e}(X_i)}}{\sum_{i=1}^n \frac{1-Z_i}{1-\hat{e}(X_i)}} . \\
\end{align}

更重要的是，结合倾向得分和 outcome 模型，可以构造双重稳健估计量：
\begin{align}
\hat{\tau}^{\mathrm{dr}} &= \hat{\tau}^{\mathrm{ht}} \label{eq:18-tau-dr} \\
&\quad - n^{-1}\sum_{i=1}^n \{Z_i - \hat{e}(X_i)\}\left\{\frac{\hat{\mu}_1(X_i)}{\hat{e}(X_i)\varepsilon_1(X_i)} + \frac{\hat{\mu}_0(X_i)\varepsilon_0(X_i)}{1-\hat{e}(X_i)}\right\} . \\
\end{align}

即，在已知 $\varepsilon_1(X_i)$ 和 $\varepsilon_0(X_i)$ 的条件下，如果倾向得分模型或 outcome 模型中的任意一个正确设定，$\hat{\tau}^{\mathrm{dr}}$ 就是 $\tau$ 的相合估计量。

\textbf{当 $\varepsilon_1(X_i) = \varepsilon_0(X_i) = 1$ 时}：上述所有估计量都退化为第 Part III 中介绍的 predictive 估计量、IPW 估计量和双重稳健估计量。

\section{R 函数与实证示例}\label{sec:18-example}

\subsection{R 函数实现}\label{sec:18-r-functions}

以下 R 函数可计算灵敏度分析的点估计：

\begin{verbatim}
OS_est_ta = function(z, y, x, out.family = gaussian,
                      truncps = c(0, 1), e1 = 1, e0 = 1) {
    # fitted propensity score
    pscore = glm(z ~ x, family = binomial)$fitted.values
    pscore = pmax(truncps[1], pmin(truncps[2], pscore))

    # fitted potential outcomes
    outcome1 = glm(y ~ x, weights = z,
                   family = out.family)$fitted.values
    outcome0 = glm(y ~ x, weights = (1 - z),
                   family = out.family)$fitted.values

    # outcome regression estimator
    ace.reg = mean(z*y) + mean((1-z)*outcome1/e1) -
              mean(z*outcome0*e0) - mean((1-z)*y)

    # IPW estimators
    w1 = pscore + (1-pscore)/e1
    w0 = pscore*e0 + (1-pscore)

    ace.ipw0 = mean(z*y*w1/pscore) - mean((1-z)*y*w0/(1-pscore))
    ace.ipw = mean(z*y*w1/pscore)/mean(z/pscore) -
              mean((1-z)*y*w0/(1-pscore))/mean((1-z)/(1-pscore))

    # doubly robust estimator
    aug = outcome1/pscore/e1 + outcome0*e0/(1-pscore)
    ace.dr = ace.ipw0 + mean((z-pscore)*aug)

    return(c(ace.reg, ace.ipw0, ace.ipw, ace.dr))
}
\end{verbatim}

\subsection{重新分析 Example 10.3}\label{sec:18-revisit-example-10-3}

沿用第 \cref{sec:10-3} 章的 NHANES 数据，研究学校午餐计划对 BMI 的影响。将灵敏度参数设为
\[
\varepsilon_1(X) = \varepsilon_0(X) \in \{1/2, 1/1.7, 1/1.5, 1/1.3, 1, 1.3, 1.5, 1.7, 2\} ,
\]
我们可以得到一列双重稳健估计。

表 \ref{tbl:18-sensitivity} 展示了灵敏度分析结果。

\begin{table}[htbp]
\caption{平均因果效应的灵敏度分析：学校午餐计划对 BMI 的影响}\label{tbl:18-sensitivity}
\[
\begin{array}{lccccccccc}
\toprule
\varepsilon & 1/2 & 1/1.7 & 1/1.5 & 1/1.3 & 1 & 1.3 & 1.5 & 1.7 & 2 \\ \midrule
1/2   & 11.62 & 10.44 & 9.40 & 8.03 & 4.96 & 0.97 & -1.69 & -4.35 & -8.34 \\
1/1.7 & 9.22  & 8.05  & 7.00 & 5.64 & 2.57 & -1.42 & -4.08 & -6.75 & -10.74 \\
1/1.5 & 7.63  & 6.45  & 5.41 & 4.05 & 0.97 & -3.02 & -5.68 & -8.34 & -12.33 \\
1/1.3 & 6.03  & 4.86  & 3.81 & 2.45 & -0.62 & -4.61 & -7.27 & -9.94 & -13.93 \\
1     & 3.64  & 2.47  & 1.42 & 0.06 & -3.01 & -7.01 & -9.67 & -12.33 & -16.32 \\
1.3   & 1.80  & 0.63  & -0.42 & -1.78 & -4.85 & -8.85 & -11.51 & -14.17 & -18.16 \\
1.5   & 0.98  & -0.19 & -1.24 & -2.60 & -5.67 & -9.66 & -12.33 & -14.99 & -18.98 \\
1.7   & 0.36  & -0.82 & -1.86 & -3.23 & -6.30 & -10.29 & -12.95 & -15.61 & -19.60 \\
2     & -0.35 & -1.52 & -2.57 & -3.93 & -7.00 & -10.99 & -13.65 & -16.32 & -20.31 \\
\bottomrule
\end{array}
\]
\end{table}

\textbf{解读}：当 $\varepsilon_1(X) > 1$ 且 $\varepsilon_0(X) > 1$ 时（即参与者倾向于有更高的 BMI），平均因果效应为负。然而，当参与者倾向于有更低的 BMI（即 $\varepsilon_1(X) < 1$ 且 $\varepsilon_0(X) < 1$）时，结论可能对灵敏度参数相当敏感。标准误差的计算留作习题 \ref{hw:18-3}。

\section{本章小结}\label{sec:18-summary}

本章围绕平均因果效应的灵敏度分析展开，核心思想是：

\begin{enumerate}
  \item \textbf{部分识别}：在没有可忽略性假设时，因果参数通常只能部分识别，即与多个值相容；
  \item \textbf{Manski 型界}：仅假设 outcomes 有界，可以导出因果效应的最坏情况界，但这些界通常过于保守；
  \item \textbf{灵敏度参数}：引入 $\varepsilon_1(X)$ 和 $\varepsilon_0(X)$ 量化未测混杂的强度；
  \item \textbf{识别公式}：在已知灵敏度参数条件下，可以给出平均因果效应的识别公式（基于 outcome regression、IPW 和双重稳健估计）；
  \item \textbf{估计量}：从识别公式出发，可以构造 predictive、projective、IPW 和双重稳健估计量；
  \item \textbf{实践应用}：灵敏度分析通过在一系列合理的 $\varepsilon$ 值上计算估计量，帮助研究者评估结论对未测混杂的稳健性。
\end{enumerate}

\section{习题}\label{sec:18-homework}

\begin{HomeworkProblem}[Proof of Theorem \ref{thm::outcome-reg-sensitivity}]\label{exr:18-1}

Prove Theorem \ref{thm::outcome-reg-sensitivity}.

\end{HomeworkProblem}

\begin{HomeworkProblem}[Proof of Theorem \ref{thm::ipw-sensitivity}]\label{exr:18-2}

Prove Theorem \ref{thm::ipw-sensitivity}.


\end{HomeworkProblem}

\begin{HomeworkProblem}[Standard errors in sensitivity analysis]\label{exr:18-3}

Chapter \ref{sec::eg-sensitivity-tau} only presents the point estimates. Report the corresponding bootstrap standard errors.




\end{HomeworkProblem}

\begin{HomeworkProblem}[Sensitivity analysis for the average causal effect on the treated units $\tau_\textsc{T}$]\label{exr:18-4}

This problem extends Chapter \ref{chapter::ATT-and-other} to allow for unmeasured confounding for estimating
\[
\tau_\textsc{T} = E\{  Y(1) - Y(0)  \mid Z=1  \}
= E(Y\mid Z=1) - E\{  Y(0)  \mid Z=1  \} .
\]
We can easily estimate $E(Y\mid Z=1)$ by the sample moment, $\hat{\mu}_{\textsc{t}1} =   \sumn Z_i Y_i / \sumn Z_i$. The only counterfactual term is $E\{  Y(0)  \mid Z=1  \}$. Therefore, we only need the sensitivity parameter $ \varepsilon_0(X) $. We have the following two identification formulas with a known  $\varepsilon_0(X) .$

\begin{Theorem}
\label{thm::att-reg-ipw}
With known $\varepsilon_0(X) $, we have
 \begin{eqnarray*}
E\{  Y(0)  \mid Z=1  \}  &=& E\left\{  Z \mu_0(X)  \varepsilon_0(X)    \right\} /e \\
&=& E\left\{  e(X)\varepsilon_0(X) \frac{ 1-Z}{1-e(X)}  Y      \right\} /e,
\end{eqnarray*}
where $e = \pr(Z=1)$
\end{Theorem}


Prove Theorem \ref{thm::att-reg-ipw}.

Remark:  Theorem \ref{thm::att-reg-ipw} motivates using $\hat{\tau}_{\textsc{t}}^* = \hat{\mu}_{\textsc{t}1}  - \hat{\mu}_{\textsc{t}0}^*$ to estimate $\tau_{\textsc{t}}$, where
 \begin{eqnarray*}
\hat{\mu}_{\textsc{t}0}^\text{reg} &=& n_1^{-1} \sumn Z_i  \varepsilon_0(X_i) \hat  \mu_0(X_i) ,   \\
\hat{\mu}_{\textsc{t}0}^\text{ht} &=& n_1^{-1} \sumn   \varepsilon_0(X_i) \hat  o(X_i) (1-Z_i) Y_i ,   \\
\hat{\mu}_{\textsc{t}0}^\text{haj} &=&  \sumn   \varepsilon_0(X_i) \hat  o(X_i) (1-Z_i) Y_i  /  \sumn   \hat  o(X_i) (1-Z_i) ,
\end{eqnarray*}
with $  \hat  o(X_i) = \hat  e(X_i) / \{1- \hat  e(X_i)\}$ being the estimated conditional odds of the treatment given the observed covariates.
Moreover, we can construct the doubly robust estimator
$\hat{\tau}_{\textsc{t}}^\text{dr} = \hat{\mu}_{\textsc{t}1}  - \hat{\mu}_{\textsc{t}0}^\text{dr}$ for $\tau_{\textsc{t}}$, where
\[
\hat{\mu}_{\textsc{t}0}^\text{dr} =
\hat{\mu}_{\textsc{t}0}^\text{ht}  - n_1^{-1} \sumn \varepsilon_0(X_i) \frac{\hat e(X_i) - Z}{1-\hat e(X_i)}  \hat \mu_0(X_i)  .
\]
\citet{luding2023flexible} provide more details.



\end{HomeworkProblem}

\begin{HomeworkProblem}[R code]\label{exr:18-5}

Implement the estimators in Problem \ref{hw::att-sentivity-analysis}.
Analyze the data used in Chapter \ref{sec::eg-sensitivity-tau}.


\end{HomeworkProblem}

\begin{HomeworkProblem}[Recommended reading]\label{exr:18-6}

 \citet{Rosenbaum::1983JRSSB} and \citet{imbens2003sensitivity} are two classic papers on sensitivity analysis which, however, involve more complicated procedures.



\end{HomeworkProblem}
